% !TEX root = ../main.tex
\chapter*{Introduction générale}
\addcontentsline{toc}{chapter}{Introduction générale}
\markboth{Introduction générale}{}

La logistique est aujourd'hui présente dans tous les secteurs économiques et constitue un
maillon stratégique du bon fonctionnement des entreprises. La présence de systèmes
informatisés de suivi des livraisons est devenue un critère essentiel pour toute organisation
souhaitant rester compétitive. Pour cette raison, les entreprises sont prêtes à investir dans des
solutions logicielles capables de centraliser et d'automatiser la gestion de leurs flux
opérationnels. En effet, compte tenu de la croissance des volumes de livraisons et de la
multiplicité des intervenants, la tâche de gérer efficacement les tournées, les chauffeurs et les
synchronisations avec les systèmes ERP s'avère de plus en plus complexe. Il devient ainsi
primordial que chaque opérateur logistique se dote d'une plateforme de gestion intégrée,
flexible et adaptée à la diversité de ses environnements techniques.

Notre contribution consiste à réaliser un projet de fin d'études en collaboration avec
l'entreprise \textbf{ASM}, basée à Sfax. L'objectif de ce projet est de concevoir et de
développer \textbf{ASM Track}, une plateforme SaaS multi-tenant de gestion et de suivi des
livraisons en temps réel. Cette solution permettra de planifier les tournées de livraison,
de suivre en temps réel la position et l'état des colis, de synchroniser automatiquement
les livraisons avec les systèmes ERP de l'entreprise cliente, et de centraliser les preuves
de livraison. Cela permettra d'éliminer les saisies manuelles, de réduire les erreurs
opérationnelles et d'offrir une meilleure visibilité à l'ensemble des parties prenantes.

Le présent rapport détaille l'avancement du projet qui s'étend sur huit chapitres, comme suit :

\textbf{Chapitre 1 :} Dans ce chapitre, nous introduisons l'entreprise d'accueil, présentons
le cadre général du projet, analysons les solutions existantes sur le marché et exposons
la méthodologie de travail adoptée.

\textbf{Chapitre 2 :} Ce chapitre présente les spécifications des besoins fonctionnels et
non fonctionnels, qui constituent la phase d'analyse du projet. Nous y établissons les
acteurs du système et les cas d'utilisation.

\textbf{Chapitre 3 :} Dans ce chapitre, nous décrivons l'architecture microservices retenue
et justifions les choix technologiques adoptés pour le backend, le frontend et le mobile.

\textbf{Chapitre 4 :} Ce chapitre décrit le développement du premier sprint : le serveur
d'authentification OAuth2 personnalisé et le cycle de vie complet des livraisons.

\textbf{Chapitre 5 :} Dans ce chapitre, nous décrivons le développement du deuxième sprint :
l'intégration ERP agnostique via le pattern Ports \& Adapters et la résilience par le
pattern Outbox transactionnel.

\textbf{Chapitre 6 :} Ce chapitre présente le troisième sprint : le constructeur de tournées,
l'intégration du moteur de routage OSRM, le suivi temps réel via WebSocket et la
surveillance des SLA.

\textbf{Chapitre 7 :} Dans ce chapitre, nous décrivons le quatrième sprint : l'application
mobile Flutter pour les chauffeurs, avec mode hors-ligne, capture de preuves de livraison
et transfert de tournée par QR code.

\textbf{Chapitre 8 :} Ce chapitre couvre le cinquième sprint : la mise en place de
l'architecture SaaS multi-tenant, le panneau super-administrateur et le déploiement
de la stack complète via Docker Compose.

Enfin, nous clôturons notre rapport par une conclusion générale exposant une synthèse du
travail réalisé et ouvrant le champ à d'autres perspectives d'évolution pour ASM Track.
